%==============================================================
%                      INTRODUCCION
%==============================================================
\chapter*{Introducción}
\addcontentsline{toc}{chapter}{Introducción}
\markboth{Introducción}{}

La gestión de la distribución, el transporte y la planificación de rutas son temas de gran importancia que han sido abordados a nivel internacional. Son ramas de la logística donde yacen numerosas oportunidades de mejora y por esta razón ha recibido especial atención motivando un intenso trabajo teórico en el desarrollo de modelos y algoritmos. En este contexto surgen los Problemas de Planificación de Rutas de Vehículos (\textit{Vehicle Routing Problems}, VRP) \parencite{colomasalazarGestionDistribucionPlanificacion2022}.

En \parencite{g.b.dantzigTruckDispatchingProblem1959} se introduce el \textit{``Truck Dispatching Problem''}, modelando cómo una flota de camiones homogéneos podría satisfacer la demanda de petróleo de varias estaciones de servicio desde un depósito central y con una distancia mínima recorrida. Cinco a\~nos después, en \parencite{clarkeSchedulingVehiclesCentral1964} se generaliza este a un problema de optimización lineal que se encuentra comúnmente en el ámbito de la logística y el transporte. El VRP es un problema NP-duro\parencite{aroraComputationalComplexityModern2009} por lo que se consideran muy difíciles de resolver. Los problemas de la vida real son considerablemente más grandes y complejos que los escenarios académicos, por lo que las heurísticas y metaheurísticas son necesarias para obtener buenos resultados. Estos algoritmos no garantizan encontrar la solución óptima, pero sí proporcionan tiempos computacionales razonables. La cantidad de métodos de solución introducidos en la literatura académica ha crecido en las últimas décadas  \parencite{braekersVehicleRoutingProblem2016}.

El Problema de Planificación de Rutas de Vehículos con Múltiples Depósitos (\textit{Multi-Depot Vehicle Routing Problems}, MDVRP) es una variante del VRP clásico que incorpora varios depósitos con una ubicación predefinida. Cada depósito cuenta con una flota de vehículos de capacidad homogénea que tiene como objetivo satisfacer la demanda de un conjunto de clientes dispersos geográficamente. Cada vehículo debe comenzar y terminar su ruta en el mismo depósito, y cada cliente debe ser visitado una única vez. Teniendo en cuenta la distancia entre la ubicación de los clientes y depósitos, el MDVRP busca asignar un cliente a cada depósito y se construye una ruta para cada vehículo dentro de un depósito para atender dichos clientes. El objetivo al crear la ruta es reducir al mínimo el costo total del recorrido \parencite{mateosDESARROLLOMETODOHIBRIDO2013}.

Durante un largo período de tiempo, los métodos heurísticos basados en la búsqueda tabú han sido los más empleados en la literatura \parencite{mateosDESARROLLOMETODOHIBRIDO2013}. Sin embargo, este problema genera una gran cantidad de combinaciones que se refleja en los altos tiempos computacionales para obtener la solución. El potencial de las técnicas de análisis de clústeres para la solución de los MDVRP ha sido reconocido por autores como \parencite{g.b.dantzigTruckDispatchingProblem1959}, \parencite{bermudezAplicacionTecnicasClusterizacion}, \parencite{barretoUsingClusteringAnalysis2007} y \parencite{mehrjerdiUsingGreedyClustering2013}. El análisis de clústeres estudia la división de objetos en grupos basado en una o varias características. Un grupo se define como una zona con muchos objetos parecidos muy cercanos entre sí, que está separada de otros grupos por zonas donde casi no hay objetos. Esta definición de grupo es bastante acorde a la agrupación de clientes por lo que se propone usar el análisis de clústeres en la solución del MDVRP, específicamente en la fase de asignación de clientes a depósitos \parencite{bermudezAplicacionTecnicasClusterizacion}. En \parencite{jainAlgorithmsClusteringData1988} se describe la técnica ''divide y vencerás'', utilizada para reducir la complejidad computacional en el empleo heurísticas para resolver el MDVRP.

Sin embargo, aún existen desafíos en cuanto a la eficiencia computacional debido a los altos tiempos y la calidad de las agrupaciones dependiendo del escenario \parencite{bermudezAplicacionTecnicasClusterizacion}. En la literatura se destacan algoritmos como \textit{Density-Based Spatial Clustering of Applications with Noise} (DBSCAN)\parencite{khanDBSCANPresentFuture2014} y \textit{K-Means}\parencite{macqueen1967some} que han sido probados y utilizados en el ámbito del MDVRP, mostrando su utilidad para la fase de asignación de clientes a depósitos  \parencite{padulohRobustOptimizingReverse2023, luo2014multi, geetha2013nested}.

En este contexto, la Universidad Tecnológica de la Habana ``José Antonio Echeverría'' (CUJAE) ha desarrollado la Biblioteca de Heurísticas de Asignación para Problemas de Planificación de Rutas de Vehículos con Múltiples Depósitos (BHAVRP). Este componente está dise\~nado para resolver la fase de asignación en MDVRP utilizando una arquitectura modular y 20 algoritmos repartidos en heurísticas clásicas de asignación y algoritmos de agrupamiento de Minería de Datos adaptados al contexto de los MDVRP \parencite{GitHubVrpsolutionsBHAVRP2026}. Esta alternativa gratuita y de código abierto supone una ventaja frente a otras herramientas que no tienen una colección de algoritmos tan extensa o son software comercial.

Esta biblioteca, en la categoría de algoritmos de agrupamiento, incluye algoritmos particionales, algoritmos jerárquicos y un algoritmo basado en densidad. En una investigación anterior \parencite{ericramosaragonIncorporacionModuloMetricas2025}, se presentó los resultados de algoritmos de agrupamiento particionales y heurísticas clásicas de asignación. Sin embargo, estos algoritmos presentaron una tendencia a asignar clientes a depósitos sin considerar adecuadamente la densidad geográfica, lo que generó clústeres con clientes muy alejados del depósito asignado. Esto se conoce en la literatura como ruido o asignaciones atípicas, lo cual puede incrementar los costos de transporte y afecta la calidad de las rutas posteriores. La literatura reconoce que los algoritmos basados en densidad, como DBSCAN, son capaces de identificar y aislar estos puntos atípicos \parencite{khanDBSCANPresentFuture2014}, ya que este tipo de algoritmos parten de la idea de que un clúster es una región de alta densidad de puntos separadas de otras regiones densas por zonas de más baja densidad.

Debido a la situación problemática antes descrita, se identifica como problema a resolver que los algoritmos de agrupamiento particionales y las heurísticas clásicas de asignación implementadas en BHAVRP generan clústeres con presencia de ruido, lo que impacta negativamente en la calidad de las rutas posteriores.

En este contexto, el objeto de estudio de este trabajo es la fase de asignación dentro de los MDVRP y la evaluación de la calidad de los algoritmos de agrupamiento. El campo de acción se enfoca en la incorporación y adaptación de algoritmos de agrupamiento basados en densidad en la fase de asignación de clientes a depósitos de BHAVRP.

Para dar solución al problema planteado se traza como objetivo general: Desarrollar una nueva versión de BHAVRP que incluya más algoritmos de agrupamiento basados en densidad.
Para dar cumplimiento al objetivo general se definieron los siguientes objetivos específicos con sus respectivas tareas:

\begin{enumerate}
  \item Identificar algoritmos de agrupamiento utilizados para resolver la fase de asignación de los problemas de planificación de rutas de vehículos con múltiples depósitos.
  \begin{itemize}
    \item Realizar búsqueda bibliográfica de las características de los problemas de planificación de rutas de vehículos con múltiples depósitos y sus aplicaciones prácticas.
    \item Realizar búsqueda bibliográfica de las propuestas de solución de MDVRP donde se realice la fase de asignación de clientes a depósitos usando como vía de resolución los algoritmos de agrupamiento.
  \end{itemize}
  \item Implementar en BHAVRP un nuevo algoritmo de agrupamiento basado en densidad.
  \begin{itemize}
    \item Analizar algoritmos de agrupamiento basados en densidad en cuanto a su funcionamiento para su incorporación en la nueva versión de la biblioteca.
    \item Seleccionar el algoritmo de agrupamiento que se van a incorporar a la biblioteca para resolver la fase de asignación del MDVRP.
    \item Implementar y adaptar el algoritmo seleccionado teniendo en cuenta las características que presentan.
  \end{itemize}
  \item Validar el funcionamiento del algoritmo mediante la ejecución  de experimentos.
  \begin{itemize}
    \item Realizar pruebas unitarias, donde se muestre que el algoritmo implementado funcione correctamente.
    \item Realizar pruebas estadísticas no paramétricas entre los resultados obtenidos por las heurísticas de construcción usando los resultados de los algoritmos de agrupamiento de BHAVRP.
    \item Realizar una prueba de Pearson y una prueba de Spearman para validar si existe una correlación positiva entre las métricas internas de BHAVRP y la calidad de las rutas obtenidas por las heurísticas de construcción.
  \end{itemize}
\end{enumerate}

El cumplimiento del objetivo general tiene como valor práctico un aporte significativo en la resolución de la fase de asignación del problema de planificación de rutas de vehículos con múltiples depósitos. Los resultados tienen valor científico y tecnológico ya que se pone a disposición de la comunidad científica un componente de software que brinda nuevos algoritmos para la solución de MDVRP.

Para alcanzar los objetivos mencionados, el trabajo se encuentra estructurado en tres capítulos, incluyendo secciones importantes como conclusiones, recomendaciones y referencias bibliográficas. El \autoref{cap:fundamentos}: \nameref{cap:fundamentos} donde se presenta el estado del arte del MDVRP, las estrategias de solución existentes, así como un análisis comparativo de bibliotecas. \autoref{cap:incorporacion}: \nameref{cap:incorporacion} describe los algoritmos de agrupamiento basados en densiad candidatos más comentados en la literatura, el algoritmo propuesto detallando el funcionamiento de este según la literatura y las adaptaciones al MDVRP que hay que realizar. Finalmente, el Capítulo 3: Análisis experimental expone varias pruebas unitarias de las funciones más relevantes del el algoritmo implementado, así como un estudio comparativo entre algoritmos de asignación en las heurísticas de construcción aplicados sobre ciertas instancias y un análisis de correlación entre las métricas de validación internas y la calidad de las rutas obtenidas.